% Pursuing Stacks (cotes 134-2 to 134-8): archive cote 134-5, site page 109 (PDF page 110), typescript page 296.
% typescript. Draft compared: ps5.tex:234-301.
% AI-transcribed by Claude Opus 5.5 (claude-opus-5-5, adaptive thinking, effort high), Message Batches API, batch ps2 (msgbatch_01NVwr1Yuo8qJf5rykeGBXPK), 2026-11-12; not yet checked by eye.
\scanpage{134-5/109}
\mspage{296}
will practically never be totally aspheric (unless we take a to be a final
object of C), hence will not be associated to a contractibility structure
-- whereas we expect that the contractibility structure of $C^\wedge$ should indu-%
ce one on $C^\wedge{}_{/a}$ . Presumably, the ``correct'' notions of induced struc-%
ture, in the case of asphericity structure and the corresponding notion
of weak equivalence, should be \xout{substant} considerably stronger than the
one I just envisioned, and correspond to the intuition of \uline{``fiberwise}
homotopy types'' over the object a (visualized as a space-like object).
A careful description of such induced structures seems to me to be\xout{ing}
needed, and \over{\unsure{a}}{t}he natural place to be the present part IV of the notes.

4) A little reflection on semi-simplicial homotopy notions (and their
analogons when \hand{$\Delta$} is replaced by a general test category\inserted{$\Delta$}) seems needed,
in order to situate the following fact: ss homotopy notions, namely for
ss objects in any category A, behave well with respect to \uline{any} functor
\[ (1) \qquad A \hand{$\longrightarrow$} \hand{$B$} , \]
without having to assume that this functor commutes with finite products,
whereas \xout{the homotopy} in the context of homotopy structures, when we
have a functor between categories endowed say with homotopy interval
structures (for instance, with contractibility structures), \xout{this}\inserted[typed]{such} func\hand{tor}
\deleted{\unsure{tion}}
\[ (2) \qquad M \hand{$\longrightarrow$} \hand{$N$} \]
behaves well with respect to homotopy notions only, it would seem, if
we assume beforehand it commutes with finite products (plus, of course,
that it transforms a given generating family of homotopy intervals \xout{\illegible{1}} of
M into homotopy intervals\xout{xx} of \hand{$N$}). In case
\[ M = \underline{Hom}(\hand{$\Delta$},\hand{$A$}) \quad , \quad \hand{$N$} = \underline{Hom}(\hand{$\Delta$},\hand{$B$}) \quad , \]
and \inserted{when} (2) comes from a functor (1) which does not commute to finite pro-%
ducts, \inserted{then} neither does (2) - and still (2) is well-behaved with respect
to \inserted{ss} homotopy notions! It should be noted of course that the semi-simpli-%
cial homotopy notions \inserted{in M} can be defined, even without assuming that in A
finite products exist, namely in situations when M does not admit finite
products - and hence, strictly speaking, the set-up of section 51 does
not apply . \inserted{All} \over{T}{t}his causes a slightly awkward feeling, which I would like
to clarify and see what's going on. I suspect it should be simple enough
to do it here and now.

First, assume that A is stable under finite products,and under
\xout{xxxxxx}\inserted[typed]{direct} sums (with small indexing set say - for what we want to do\inserted{with $\Delta$,}
finite direct sums even would be enough). We'll use sums only with sum\hand{-}\inserted{mands}

\margin{109}
